\uuid{DqQZ}
\titre{Changement de variables}
\auteur{hanani ; mzari}
\organisation{exo7}
\licence{CC BY-NC-SA 4.0}
\datecreate{2020-01-20}
\niveau{L1}
\module{Analyse}
\chapitre{Calcul d'intégrales}
\sousChapitre{Primitives diverses}
\difficulte{2}
\theme{Primitives, changement de variables}
\exo7qcmid{lille-2.primitives.019}

\contenu{
\begin{question}
\qclassification{127.11}
\qtags{motcle=changement de variables}
Parmi les affirmations suivantes, cocher celles qui sont vraies :
\begin{answers}
\bad{$\displaystyle \int \frac{x\mathrm{d}x}{1+x^2}=\arctan (x)+k$, où $k\in \Rr$.}
\bad{$\displaystyle \int x\mathrm{e}^{(1+x^2)}\mathrm{d}x=\frac{x^2}{2}\mathrm{e}^{(1+x^2)}+k$, où $k\in \Rr$.}
\good{$\displaystyle \int \frac{x\mathrm{d}x}{1+x^2}=\ln\sqrt{1+x^2}+k$, où $k\in \Rr$.}
\good{$\displaystyle \int x\mathrm{e}^{(1+x^2)}\mathrm{d}x=\frac{1}{2}\mathrm{e}^{(1+x^2)}+k$, où $k\in \Rr$.}
\end{answers}
\begin{explanations}
Avec $u=1+x^2$, on a : $\mathrm{d}u=2x\mathrm{d}x$. D'où
$$\int \frac{x\mathrm{d}x}{1+x^2}=\frac{1}{2}\int \frac{\mathrm{d}u}{u}=\frac{1}{2}\ln (u)+k=\ln\sqrt{1+x^2}+k,\; k\in \Rr,$$
et
$$\int x\mathrm{e}^{(1+x^2)}\mathrm{d}x=\frac{1}{2}\int \mathrm{e}^u\mathrm{d}u=\frac{1}{2}\mathrm{e}^u+k=\frac{1}{2}\mathrm{e}^{(1+x^2)}+k,\; k\in \Rr.$$
\end{explanations}
\end{question}
}
